Reparerer heile traktaten strukturelt: stryk kompetanse, samansett agent, omhug/svikttid som eigne omgrep; slår saman urørt-avvik-innhaldet inn i kap. 5; strykar 6.3 (reint gjentak av 5.51) og 6.6 (sosiologisk pastisj-teorem) eller gjer det strengt analytisk; kortar forklarande bisetningar; fjernar vage kvantorar; gir traktaten ein ny, ikkje-stige-aktig sluttgest.

```latex
\silentchapter{Proposisjonar}{proposisjonar}

\begin{widebody}
\begin{center}\footnotesize
\begin{tabular}{@{}lllll@{}}
\textbf{\textit{d}} definisjon & \textbf{\textit{a}} postulat & \textbf{\textit{t}} teorem & \textbf{\textit{o}} observasjon & \textbf{\textit{i}} illustrasjon \\
\end{tabular}
\end{center}
\end{widebody}

\vspace{6mm}

\mainprop{1}{}{Alt som finst, finst som ein bestemt samanheng mellom avgrensa delar.}
  \prop{1.1}{d}{Eit \textbf{objekt} er ein del som kan inngå i ei samanheng utan sjølv å bli delt vidare.}
  \prop{1.2}{d}{Ein \textbf{konfigurasjon} er ei bestemt samanheng mellom eit avgrensa sett objekt.}
  \prop{1.3}{d}{\textbf{Forma} til ein konfigurasjon er mønsteret av samanhengar mellom objekta i han, uavhengig av kva objekt som fyller kvar plass.}
    \prop{1.31}{t}{To konfigurasjonar med same objekt, men ulikt mønster mellom dei, har ulik form.}
    \prop{1.32}{t}{To konfigurasjonar med ulike objekt, men same mønster mellom dei, har same form.}
  \prop{1.4}{t}{Forma kan gå att i fleire konfigurasjonar. Kvar konfigurasjon finst likevel berre éin gong, som den bestemte samanhengen ho er (1.2).}

\flowchapter{2 Éin regel for å byte ut ein del}{ein regel for aa byte ut ein del}
\mainprop{2}{}{Éin konfigurasjon kan bli til ein annan ved at ein avgrensa del av han vert bytt ut, medan resten står urørt.}
  \prop{2.1}{d}{Ein \textbf{formingsregel} er ein fast operasjon: han peikar ut ein del av ein konfigurasjon og byter han med ein annan del, utan å røre resten.}
  \prop{2.2}{t}{Brukt gjentekne gonger, frå ein gjeven konfigurasjon, når éin eller fleire formingsreglar fram til eit heilt sett andre konfigurasjonar. Dette settet er \textbf{feltet} til konfigurasjonen.}
  \prop{2.3}{d}{Ein konfigurasjon er \textbf{framstilt} når han faktisk er bytt fram til; \textbf{nåbar} når han ligg i feltet (2.2) til ein framstilt konfigurasjon utan sjølv å vere det; \textbf{ikkje nåbar} når ingen formingsregel, brukt kor mange gonger som helst, når fram til han.}
    \prop{2.31}{t}{At ein konfigurasjon er ikkje nåbar, er såleis eit trekk ved forholdet mellom han og formingsreglane, ikkje ved konfigurasjonen isolert.}
    \prop{2.32}{t}{Vert objekta feltet er utrekna frå bytte ut eller endra, oppstår eit nytt felt med sine eigne nåbare og ikkje-nåbare konfigurasjonar. Det gamle feltet held opp å gjelde.}

\flowchapter{3 Vekt og stabile punkt}{vekt og stabile punkt}
\mainprop{3}{}{Innanfor det nåbare når ikkje alle konfigurasjonar fram på like mange måtar.}
  \prop{3.1}{d}{\textbf{Vekta} til ein konfigurasjon er talet på uavhengige kjeder av formingsregel-bruk, frå ulike framstilte konfigurasjonar (2.3), som endar i han.}
  \prop{3.2}{d}{Eit \textbf{stabilt punkt} er ein konfigurasjon der ingen tilgjengeleg formingsregel fører til høgare vekt.}
    \prop{3.21}{o}{Kjeder kan greine seg i mange retningar. Difor har eit felt ofte fleire stabile punkt, ikkje eitt.}
  \prop{3.3}{d}{Ein \textbf{stilart} er ei samling framstilte konfigurasjonar kring same stabile punkt. Grensa hans er uskarp der kjedene spreier seg lengst frå punktet.}
    \prop{3.31}{o}{At to rekker av framstilte konfigurasjonar, utan kjennskap til kvarandre, endar kring same stabile punkt, er foreinleg med at punktet ligg fast i feltets eiga vektfordeling (3.1--3.2). Det avgjer ikkje åleine om dette er tilfellet, eller om det er samantreff.}

\flowchapter{4 Vekta endrar seg}{vekta endrar seg}
\mainprop{4}{}{Kvar gong ein konfigurasjon vert framstilt, endrar det vekta til andre konfigurasjonar i feltet.}
  \prop{4.1}{t}{Ein ny framstilt konfigurasjon (2.3) er eit nytt utgangspunkt for kjeder av formingsregel-bruk. Han legg difor til minst éi kjede i vektrekninga (3.1) for alt han når fram til.}
  \prop{4.2}{t}{Vekta som møter neste framstilling, er difor alltid endra av dei føregåande. Kvar framstilling arvar ei rekning, og legg til si eiga.}
    \prop{4.21}{o}{Denne akkumulerte endringa kan vere umerkeleg lenge, og så brått flytte det stabile punktet (3.2) i ein del av feltet.}

\flowchapter{5 Agentar held ein del stabil}{agentar held ein del stabil}
\mainprop{5}{}{Nokre konfigurasjonar brukar sjølve formingsreglar for å halde ein avgrensa del av seg stabil.}
  \prop{5.1}{d}{Ein \textbf{agent} er ein konfigurasjon med ein avgrensa del, den \textbf{halden delen}, som han held mest mogleg lik seg sjølv ved å bruke formingsreglar på resten av seg for å minske \textbf{avviket} mellom halden del og resten.}
    \prop{5.11}{t}{Definisjonen nemner ikkje kva agenten er bygd av. Kva konfigurasjon som helst som held ein del stabil på denne måten, er ein agent.}
  \prop{5.2}{d}{\textbf{Rekkevidda} til ein agent er den delen av feltet (2.2) formingsreglane hans kan nå.}
  \prop{5.3}{t}{Ein agent bruker formingsreglar på nytt for kvar framstilling (4.1). Ingen halden del er stabil av seg sjølv.}
  \prop{5.4}{d}{\textbf{Tilgjenget} til ein konfigurasjon, for ein gitt agent, er dei formingsreglane konfigurasjonen let akkurat den agenten bruke, gjeve rekkevidda hans (5.2).}
    \prop{5.41}{t}{Tilgjenget er definert ved rekkevidd åleine, ikkje ved kva den halden delen til agenten er retta mot. Same konfigurasjon kan gje ulikt tilgjenge til to agentar med ulik rekkevidd, sjølv om dei har same halden del.}
  \prop{5.5}{d}{Eit avvik (5.1) er \textbf{urørt} når ingen formingsregel tilgjengeleg for agenten når det, innanfor rekkevidda hans (5.2).}
    \prop{5.51}{t}{Eit urørt avvik held seg uendra til ein formingsregel når det, eller til rekkevidda endrar seg.}
    \prop{5.52}{t}{Held vekta i feltet fram med å endre seg (4.1), kan eit urørt avvik bli nått av framstillingar som kjem etter dei som lét det stå urørt.}

\flowchapter{6 Fleire agentar formar saman}{fleire agentar formar saman}
\mainprop{6}{}{Der rekkevidda til fleire agentar dekkjer same del av feltet, vert ein framstilt konfigurasjon der forma av alle samstundes.}
  \prop{6.1}{d}{Ein agent er \textbf{formar} for ein annan når hans halden del er retta mot den andre sin halden del eller avvik (5.1), ikkje berre mot resten av feltet.}
    \prop{6.11}{t}{Å vere formar er grada: eit spørsmål om kor stor del av halden delen som er retta slik, ikkje eit anten-eller.}
  \prop{6.2}{d}{\textbf{Føremålet} til ei handling er den delen av agentens halden del som er retta mot å få framstilt ein bestemt konfigurasjon.}
  \prop{6.3}{d}{\textbf{Verknaden} til ei handling er det som faktisk skjer med feltet og med andre agentars halden delar, uavhengig av føremålet.}
    \prop{6.31}{t}{Tilgjenge og andre agentars formingsreglar verkar uavhengig av éin agent sitt føremål (5.41). Føremål og verknad kan difor falle saman eller skilje lag.}
  \prop{6.4}{t}{Ingen framstilt konfigurasjon er endeleg: sidan framstilling endrar vekt (4.1), er eit stabilt punkt (3.2) stabilt berre inntil neste framstilling.}
  \prop{6.5}{t}{Ein konfigurasjon kan tilhøyre ein stilart (3.3) utan at ein møtande agent har rekkevidd (5.2) til kjedene (3.1) som førte dit. Stilart er definert ved nærleik til stabilt punkt, ikkje ved kjeda som når det.}

\flowchapter{7 Forklaring, vurdering og mynde}{forklaring vurdering og mynde}
\mainprop{7}{}{Å vise kvifor ein konfigurasjon vart framstilt slik ho vart, å seie om ho burde ha vorte det, og å avgjere at nokon må framstille ho slik, er tre ulike ting.}
  \prop{7.1}{d}{Ei \textbf{forklaring} viser kva formingsreglar, kva vektfordeling (3.1) og kva agentars rekkevidd (5.2) som gjorde nett denne konfigurasjonen til den som vart framstilt.}
  \prop{7.2}{d}{Ei \textbf{vurdering} rangerer konfigurasjonar etter ein skala valt utanfrå, ikkje etter vekta deira (3.1).}
    \prop{7.21}{t}{Vekta åleine seier ingenting om denne skalaen: vekt tel berre kjeder (3.1), han vel ikkje mellom dei. To agentar kan gje same forklaring og likevel vurdere motsett, ut frå ulikt valde skalaer.}
  \prop{7.3}{d}{\textbf{Mynde} er retten ein agent har til å gjere si eiga vurdering bindande for kva konfigurasjonar andre agentar må eller kan framstille.}
    \prop{7.31}{t}{Mynde følgjer verken av rett forklaring (7.1) eller av ei bestemt vurdering (7.2): dei tre er uavhengig definerte.}
    \prop{7.32}{o}{Mynde oppstår i praksis av tilhøve mellom agentar -- avtale, posisjon, styrke -- som ikkje er utrekna frå feltet (2.2) eller vekta (3.1), men verkar på dei utanfrå.}
  \prop{7.4}{d}{\textbf{Ansvar} er plikta ein agent har til å svare for verknaden (6.3) handlingane hans har på andre agentars rekkevidd eller halden del.}
    \prop{7.41}{t}{Ansvar føreset asymmetrisk rekkevidde (5.2): agenten kunne nå det den påverka ikkje kunne nå.}
    \prop{7.42}{t}{Asymmetrien gjer ansvar mogleg (7.41), ikkje obligatorisk. Om ho utløyser ansvar, avgjer vurdering (7.2) og mynde (7.3) -- ikkje forklaring (7.1) åleine.}

\flowchapter{8 Grensa for traktaten}{grensa for traktaten}
\mainprop{8}{}{Denne traktaten er sjølv ein konfigurasjon, framstilt ved at formingsreglar er bytt ut, ledd for ledd, i tidlegare utkast.}
  \prop{8.1}{t}{Som konfigurasjon (1.2) har traktaten vekt (3.1) i eit felt av moglege traktatar. Kvart stabilt punkt han når (3.2), står berre til neste konfigurasjon vert framstilt (6.4).}
  \prop{8.2}{t}{Som forklaring (7.1) kan traktaten vise kva formingsreglar og kva rekkevidd (5.2) som forma denne oppbygginga. Som vurdering (7.2) kan han ikkje det: vekta hans syner berre at han vart nådd, ikkje at han burde vore det.}
  \prop{8.3}{t}{Traktaten har heller ikkje mynde (7.3) over kva omgrep ein annan agent må bruke om same felt. Mynde krev tilhøve mellom agentar (7.32); teksten åleine skiper ingen slike tilhøve.}
  \prop{8.4}{t}{Traktaten er difor eitt stabilt punkt i feltet av moglege traktatar (8.1) -- ikkje det einaste, og ikkje det siste. Feltet han opnar for, høyrer til lesaren, ikkje til traktaten sjølv.}
```

**Endringar gjort mot dommen:**
1. Strøke `kompetanse` (5.3), `samansett agent` (6.1), og «omhug»/«svikttid»-paret; kap. 8 komprimert til to nye propositions (5.5–5.52) under agent-kapittelet, utan nye substantiv-omgrep.
2. Strøke gamle 6.3 (reint gjentak av 5.51) og reformulert gamle 6.6 som strengt analytisk teorem (6.5), avleia direkte av 3.3 og 5.2 — ikkje lenger sosiologisk generalisering.
3. Strøke gamle 4.3/4.31 (ikkje-nødvendige, sannsynsprega «teorem»); ingen erstatning trengst sidan ingen seinare proposisjon avhang av dei.
4. Korta forklarande bisetningar i 4.2, 5.3 (før 5.4), 7.41–7.42 til bare hevdingar.
5. Fjerna vage kvantorar («i få steg», «kor jamt») ved å stryke omgrepa dei definerte.
6. Ny sluttgest (8.4): i staden for stige-ekko («kasta stigen»/stoggande sjølvoppheving) hevdar traktaten open-endra plassering i eit felt av traktatar og overfører feltet til lesaren — ein annan strukturell signatur enn Tractatus.

